% ============================================================================
% WORKBOOK — Section 6.6: Angle Relationships
% CCSS 7.G.B.5
% Practice-focused reinforcement for the study guide topic
% ============================================================================
\section{Angle Relationships}
\topicTitle{Angle Relationships}

\begin{quickReview}{Angle Pairs}
\begin{itemize}[leftmargin=*]
    \item \textbf{Complementary angles} add to $90°$. \quad Example: $35° + 55° = 90°$.
    \item \textbf{Supplementary angles} add to $180°$. \quad Example: $110° + 70° = 180°$.
    \item \textbf{Vertical angles} are formed by two intersecting lines and are always \textbf{equal}.
    \item \textbf{Adjacent angles} share a side and a vertex.
\end{itemize}

\medskip
\textbf{To find an unknown angle:} set up an equation using the relationship, then solve for the variable.

\smallskip
\textbf{Example:} Two supplementary angles are $x$ and $3x$. Then $x + 3x = 180$, so $4x = 180$ and $x = 45°$. The angles are $45°$ and $135°$.
\end{quickReview}

% ── Warm-Up ──────────────────────────────────────────────────────────────────
\begin{practiceBox}[funGreen]{\faSun~Warm-Up}

\practiceHeader[funGreen]{Find the Missing Angle}

\begin{multicols}{2}
\prob Complement of $25°$: \answerBlank[2cm]
\answerExplain{$65°$}{Complementary angles sum to $90°$: $90° - 25° = 65°$.}

\prob Supplement of $110°$: \answerBlank[2cm]
\answerExplain{$70°$}{Supplementary angles sum to $180°$: $180° - 110° = 70°$.}

\prob Complement of $48°$: \answerBlank[2cm]
\answerExplain{$42°$}{$90° - 48° = 42°$.}

\prob Supplement of $73°$: \answerBlank[2cm]
\answerExplain{$107°$}{$180° - 73° = 107°$.}

\prob Complement of $62°$: \answerBlank[2cm]
\answerExplain{$28°$}{$90° - 62° = 28°$.}

\prob Supplement of $155°$: \answerBlank[2cm]
\answerExplain{$25°$}{$180° - 155° = 25°$.}
\end{multicols}
\end{practiceBox}

% ── Practice A: Identify the Relationship ────────────────────────────────────
\begin{practiceBox}{Name That Angle Pair}

\practiceHeader{Write \textbf{complementary}, \textbf{supplementary}, or \textbf{vertical} for each pair.}

\begin{multicols}{2}
\prob $30°$ and $60°$ \answerBlank[3cm]
\answerExplain{Complementary}{$30° + 60° = 90°$. Angles that sum to $90°$ are complementary.}

\prob $95°$ and $85°$ \answerBlank[3cm]
\answerExplain{Supplementary}{$95° + 85° = 180°$. Angles that sum to $180°$ are supplementary.}

\prob Two equal angles formed by intersecting lines, each $40°$. \answerBlank[3cm]
\answerExplain{Vertical}{When two lines intersect, opposite angles are vertical angles and are always equal.}

\prob $72°$ and $18°$ \answerBlank[3cm]
\answerExplain{Complementary}{$72° + 18° = 90°$. These are complementary angles.}

\prob $135°$ and $45°$ \answerBlank[3cm]
\answerExplain{Supplementary}{$135° + 45° = 180°$. These are supplementary angles.}

\prob $53°$ and $37°$ \answerBlank[3cm]
\answerExplain{Complementary}{$53° + 37° = 90°$. These are complementary angles.}
\end{multicols}
\end{practiceBox}

% ── Practice B: Solve Equations for Angles ───────────────────────────────────
\begin{practiceBox}[funTeal]{Solving for Unknown Angles}

\practiceHeader[funTeal]{Set up and solve an equation to find the unknown angle(s).}

\prob Two complementary angles are $x$ and $2x$. Find both angles. \answerBlank[3cm]
\answerExplain{$30°$ and $60°$}{$x + 2x = 90°$. Combine: $3x = 90$. Divide: $x = 30$. The angles are $30°$ and $2(30) = 60°$. Check: $30 + 60 = 90°$ \checkmark.}

\prob Two supplementary angles are $(x + 20)°$ and $(3x)°$. Find both angles. \answerBlank[3cm]
\answerExplain{$60°$ and $120°$}{$x + 20 + 3x = 180$. Combine: $4x + 20 = 180$. Subtract $20$: $4x = 160$. Divide: $x = 40$. Angles: $40 + 20 = 60°$ and $3(40) = 120°$. Check: $60 + 120 = 180°$ \checkmark.}

\prob Two vertical angles are $(5x - 10)°$ and $(3x + 20)°$. Find $x$ and the angle measure. \answerBlank[3cm]
\answerExplain{$x = 15$; angles are $65°$}{Vertical angles are equal: $5x - 10 = 3x + 20$. Subtract $3x$: $2x - 10 = 20$. Add $10$: $2x = 30$. Divide: $x = 15$. Angle: $5(15) - 10 = 65°$. Check: $3(15) + 20 = 65°$ \checkmark.}

\prob Two supplementary angles are in the ratio $3:7$. Find both angles. \answerBlank[3cm]
\answerExplain{$54°$ and $126°$}{Let the angles be $3x$ and $7x$. Sum: $3x + 7x = 180$, so $10x = 180$ and $x = 18$. Angles: $3(18) = 54°$ and $7(18) = 126°$. Check: $54 + 126 = 180°$ \checkmark.}

\prob Two complementary angles are $(4x + 5)°$ and $(x + 10)°$. Find both angles. \answerBlank[3cm]
\answerExplain{$65°$ and $25°$}{$4x + 5 + x + 10 = 90$. Combine: $5x + 15 = 90$. Subtract $15$: $5x = 75$. Divide: $x = 15$. Angles: $4(15) + 5 = 65°$ and $15 + 10 = 25°$. Check: $65 + 25 = 90°$ \checkmark.}
\end{practiceBox}

% ── Practice C: True or False ────────────────────────────────────────────────
\begin{practiceBox}[funPurple]{True or False?}

\trueOrFalse{Two vertical angles can be supplementary.}
\answerExplain{True}{Vertical angles are equal. If both equal $90°$, they are also supplementary ($90° + 90° = 180°$). This happens when two lines are perpendicular.}

\trueOrFalse{A pair of complementary angles can include an obtuse angle.}
\answerExplain{False}{Complementary angles sum to $90°$. If one angle is obtuse (greater than $90°$), the other would need to be negative, which is impossible.}

\trueOrFalse{If two angles are supplementary and one is $90°$, the other is also $90°$.}
\answerExplain{True}{$180° - 90° = 90°$. Both angles are right angles.}

\trueOrFalse{Vertical angles are always congruent.}
\answerExplain{True}{Vertical angles are formed by two intersecting lines, and the opposite angles are always equal (congruent).}
\end{practiceBox}

% ── Word Problems ────────────────────────────────────────────────────────────
\begin{practiceBox}[funBlue]{\faGlobe~Word Problems}

\wordProblem{Two streets intersect. One of the four angles formed is $65°$. What are the measures of the other three angles?}{degrees}
\answerExplain{$65°$, $115°$, $115°$}{The vertical angle is also $65°$. Each of the two supplementary angles is $180° - 65° = 115°$.}

\wordProblem{The corner of a bookshelf meets a wall, forming two complementary angles. One angle is $15°$ more than the other. Find both angles.}{degrees}
\answerExplain{$37.5°$ and $52.5°$}{Let the smaller angle be $x$. Then $x + (x + 15) = 90$. Combine: $2x + 15 = 90$. Subtract $15$: $2x = 75$. Divide: $x = 37.5$. The angles are $37.5°$ and $52.5°$.}

\wordProblem{An angle is three times its supplement. Find the angle.}{degrees}
\answerExplain{$135°$}{Let the angle be $x$ and its supplement be $180 - x$. Then $x = 3(180 - x)$. Distribute: $x = 540 - 3x$. Add $3x$: $4x = 540$. Divide: $x = 135°$. Its supplement is $180 - 135 = 45°$. Check: $135 = 3 \times 45$ \checkmark.}
\end{practiceBox}

% ── Challenge ────────────────────────────────────────────────────────────────
\begin{challengeBox}
\prob Two intersecting lines form four angles. One angle is $(2x + 30)°$. The adjacent angle is $(3x - 10)°$. Find all four angles.
\answerExplain{$94°$, $86°$, $94°$, $86°$}{Adjacent angles at intersecting lines are supplementary: $(2x + 30) + (3x - 10) = 180$. Combine: $5x + 20 = 180$. Subtract $20$: $5x = 160$. Divide: $x = 32$. Angles: $2(32) + 30 = 94°$ and $3(32) - 10 = 86°$. The four angles are $94°$, $86°$, $94°$, $86°$ (vertical pairs). Check: $94 + 86 = 180$ \checkmark.}

\prob An angle is equal to its complement. What is the angle?
\answerExplain{$45°$}{Let the angle be $x$. Its complement is $90 - x$. Set equal: $x = 90 - x$. Add $x$: $2x = 90$. Divide: $x = 45°$. Check: the complement of $45°$ is $90 - 45 = 45°$ \checkmark.}
\end{challengeBox}

\encouragement{Outstanding angle work! You have mastered complementary, supplementary, and vertical angles!}
